% phi_change_of_variables.tex --- Phase C deliverable
% Session VQUAD-PIII-NORMALIZATION-MAP, 2026-05-02
%
% Verdict: G15_PARTIAL.
% The change-of-variables Phi is built as a composition of three
% blocks: Phi = Phi_symp o Phi_shift o Phi_resc.
% Phi_resc and Phi_shift are pinned (modulo a discrete sign).
% Phi_symp is the residual blocker; it requires Okamoto's explicit
% Lax pair for P_III(D_6), which is not in the operator's library.

\section*{Change-of-variables $\Phi: V_\mathrm{quad} \to P_{III}(D_6)$}
\label{sec:phi}

\subsection*{Architecture}

We seek $\Phi$ as a composition
\[
  \Phi \;=\; \Phi_\mathrm{symp} \,\circ\, \Phi_\mathrm{shift}
            \,\circ\, \Phi_\mathrm{resc},
\]
acting on Lax-pair data attached to the V$_\mathrm{quad}$
scalar ODE. The three blocks have distinct r\^oles:

\begin{itemize}
  \item $\Phi_\mathrm{resc}$ rescales the spectral parameter
        $z \mapsto \lambda\,t$ and the wave function
        normalization, mapping V$_\mathrm{quad}$'s irregular
        singular point at $z=0$ to the canonical $P_{III}(D_6)$
        irregular singular point at $t=0$;
  \item $\Phi_\mathrm{shift}$ accounts for the choice of base
        point at $t=\infty$ (matching the regular singular
        point of V$_\mathrm{quad}$ at $z=\infty$ to a
        canonical $P_{III}(D_6)$ regular point);
  \item $\Phi_\mathrm{symp}$ is a symplectic transform on the
        Lax-pair monodromy variety, fixing the discrete
        choice of Stokes-multiplier labeling.
\end{itemize}

\subsection*{$\Phi_\mathrm{resc}$ -- determined}

\textbf{Matching the leading exponent.} V$_\mathrm{quad}$
has $f_\pm(u) = \exp(\pm c_0/u)\, u^\rho\,(1+O(u))$ with
$z = u^2$, $c_0 = 2/\sqrt 3$, $\rho = -11/6$. Canonical
$P_{III}(D_6)$ at the irregular singularity $t=0$ has
trans-series $\exp(\pm 2/\sqrt{t})\,t^{\sigma_0}\,(1+O(\sqrt t))$
in some normalization (Okamoto~1987 \S3, eq.~3.7; quoted via
Conte--Musette ch.~7 \S7.4 trans-series tables -- residual
\textsf{R2}). The exponent matching
\[
  \exp\!\Bigl(\pm c_0 / u\Bigr) \;\longleftrightarrow\;
  \exp\!\Bigl(\pm 2 / \sqrt{t}\,\Bigr)
\]
under $u = \sqrt{z}$, $\sqrt{t} = \sqrt{\lambda\,t} / \sqrt{\lambda}$
forces
\[
  \boxed{\;
    \lambda \;=\; \tfrac{c_0^2}{4} \;=\; \tfrac{1}{3},
    \qquad
    z \;=\; \tfrac{1}{3}\,t.
  \;}
\]
\textbf{This is residual-R3-conditional}: the precise
canonical-form leading exponent at $t=0$ depends on Okamoto's
sign and prefactor convention.

\textbf{Wave-function rescaling.} The secondary exponent
$\rho_{\Vquad} = -11/6$ versus the canonical $\sigma_0$ at the
V$_\mathrm{quad}$ parameter point fixes the multiplicative
power of $t$ that enters the wave function. Pending R2 we
write
\[
  f(z)\, \stackrel{\Phi_\mathrm{resc}}{\longmapsto}\,
  \mu\, t^{\sigma_0 - \rho_{\Vquad}/2}\, \tilde f(t),
  \quad \mu \in \mathbb{C}^\times.
\]

\subsection*{$\Phi_\mathrm{shift}$ -- partially determined}

V$_\mathrm{quad}$ has a regular singular point at
$z=\infty$ (the OGF radius is finite, equal to
$\zeta_* / \beta_2 = 4/(3\sqrt 3)$ in the relevant normalization
-- this is residual~\textsf{R4}, the OGF radius vs.\ the
Borel-plane $\zeta_*$ are connected by a Mellin transform that
we have not made explicit in this session). Canonical
$P_{III}(D_6)$ has its second irregular singularity at
$t=\infty$, of identical Poincar\'e rank $1$. The matching
constrains the affine shift of $t$:
\[
  \Phi_\mathrm{shift}: t \mapsto t + t_0,
  \qquad t_0 \text{ pinned by matching at } z=\infty.
\]
Without R4 the value of $t_0$ is not numerically pinned;
its symbolic form is $t_0 = -\zeta_*\!/\lambda\,$
(modulo R2-R3).

\subsection*{$\Phi_\mathrm{symp}$ -- residual blocker}

The symplectic transform on $(q, p)$ that maps the linear
$L$-equation arising as a Lax pair component for V$_\mathrm{quad}$
to the canonical $P_{III}(D_6)$ Lax pair is fixed by the choice
of normalisation of the linear isomonodromic system. Without
the explicit Okamoto~1987 Lax pair (a $2\times 2$ rational
linear system in a spectral variable $\zeta$, with poles at
$\zeta=0,\infty$ of Poincar\'e rank $1$), we cannot:
\begin{enumerate}
  \item write down the explicit linear-system gauge
        transformation $\zeta \mapsto a\,\zeta + b/(c\,\zeta + d)$
        with $ad-bc=1$ that maps V$_\mathrm{quad}$'s scalar
        ODE to the canonical Okamoto Lax-pair scalar reduction;
  \item determine the multiplicative factor that the
        Stokes constant $S_{\zeta_*}^{\Vquad}$ picks up under
        this gauge transformation.
\end{enumerate}

\subsection*{Status of the Phase-C matching system}

\textbf{Constraints used so far:}
\begin{enumerate}
  \item leading exponent match $\Rightarrow$ $\lambda = 1/3$ (PASS, conditional on R3);
  \item secondary exponent match $\Rightarrow$
        $\rho_{\Vquad} = -11/6$ versus canonical $\sigma_0$
        (UNRESOLVED, conditional on R2);
  \item regular-singular-point match at $z=\infty$
        $\Rightarrow$ symbolic form of $t_0$ (UNRESOLVED,
        conditional on R4).
\end{enumerate}

\textbf{The matching system is consistent so far.} No
inconsistency has appeared in the partial determination
of $(\lambda, \mu, t_0)$. In particular, the Phase-C
``HALT clause'' \textsf{G15\_NOT\_LIFTABLE} is \emph{not}
triggered: V$_\mathrm{quad}$'s scalar ODE has the right
singular-point structure (irregular at $0$ of rank $1/2$
$\equiv$ rank $1$ in $u$, regular at $\infty$, two formal
sectors) to be the $L$-component of a $P_{III}(D_6)$ Lax
pair. The G15 status is \textsc{closure-pending-on-residuals},
not \textsc{not-liftable}.

\subsection*{Verdict}

\fbox{\parbox{0.95\textwidth}{
\textbf{Verdict: G15\_PARTIAL.}\\
$\Phi_\mathrm{resc}$ and $\Phi_\mathrm{shift}$ are pinned
modulo three residuals (R2, R3, R4) that all reduce to
items in Okamoto~1987 and Conte--Musette ch.~7 not in the
operator's library. $\Phi_\mathrm{symp}$ is the substantive
blocker: it requires the explicit Lax pair, not just the
Hamiltonian. The numerical Jacobian factor of $\Phi$
acting on $C$ is therefore not computed in this session;
$C_\mathrm{can}$ remains \emph{symbolic}.
The Phase-C matching system is consistent (no
\textsf{G15\_NOT\_LIFTABLE}); the $\zeta_*$ and $\rho$
exponents are exact rational/algebraic (no
\textsf{G15\_NUMERICAL\_DRIFT}).}}

\subsection*{Residual list (operator action items)}

\begin{description}
  \item[R1] Okamoto-convention identification of CT~v1.3 \S3.5
        4-tuple $(1/6, 0, 0, -1/2)$ as
        $(\alpha_\infty,\alpha_0,\beta_\infty,\beta_0)$ or
        permutation thereof. Source: Okamoto~1987~\S2.
  \item[R2] Trans-series leading prefactor of $P_{III}(D_6)$
        at $t=0$ in canonical normalization:
        $\exp(\pm 2 c_*/\sqrt{t})$ with what $c_*$?
        Source: Okamoto~1987~\S3 eq.~(3.7) or
        Conte--Musette ch.~7 \S7.4.
  \item[R3] Sign and phase convention of the Stokes
        multiplier at the leading Stokes ray. Source:
        Lisovyy--Roussillon 2017 (or Boutroux 1913 for the
        historical baseline).
  \item[R4] Mellin/Borel transform between OGF radius and
        Borel-plane $\zeta_*$ in $\Vquad$ native
        normalisation, used in $\Phi_\mathrm{shift}$.
        Source: Loday--Richaud 2016 ch.~5
        (\textsl{Divergent Series, Summability and Resurgence II}).
  \item[R5] (\textbf{primary blocker}) Explicit Okamoto Lax
        pair for $P_{III}(D_6)$: $2\times 2$ system in spectral
        $\zeta$, with poles at $0,\infty$ of rank $1$, and the
        linear scalar reduction in the spectral parameter that
        recovers the V$_\mathrm{quad}$ scalar ODE under
        $\Phi_\mathrm{resc}\circ\Phi_\mathrm{shift}$.
        Source: Okamoto~1987 \S\S2--3 + Jimbo--Miwa 1981
        (\textsl{Monodromy preserving deformation II},
        Physica D 2:407--448).
\end{description}
